
\section{Introduction}
Deceptive practices have been been a documented problem and research topic on Web Studies since the Web 1.0 era \cite{wwwCloak2006}

\textbf{RQ1} What forms of URL cloaking can be detected internally within the Ad Library? Which mechanisms are exploited and how prevalent are they?


\textbf{RQ2} What is the distribution of domains embedded in adverts? When requested, what kinds of redirection cloaking can be observed?



\textbf{RQ3} What features further differ cloaked and non-cloaked adverts?



This studies contributions to literature are:
\begin{itemize}
    \item A first of its kind analysis of advert-embedded URLs and cloaking patterns employed in Meta's Advertising Ecossystem, leading to the uncovering of unintended, malicious abuse of certain Ad Distribution features with the goal of avoiding moderation and deceiving consumers.
    \item A unified definition of cloaking over more than 20 years of dedicated research, leading to a highlighting of enduring, ressurging and novel deceitful practices and exposing their relation to adjacent research areas.
    \item A discussion on low-resource moderation strategies that could effectively reduce deceitful patterns. 
\end{itemize}
